\documentclass[10pt,a4paper]{amsart}
\usepackage[margin=19mm]{geometry}
\usepackage{amsmath,amssymb,mathtools}
\usepackage[scr=rsfs,cal=euler]{mathalfa}
\usepackage{xfrac}
\usepackage[unicode,hidelinks,bookmarks=false]{hyperref}
\newcommand{\ie}{i.e.\ }
\newcommand{\cf}{cf.\ }
\newcommand{\resp}{respectively\ }
\newcommand{\Subsection}{\subsection}
\providecommand{\nobreakdash}{\nobreak\hbox{-}}
\setlength{\emergencystretch}{2em}
\multlinegap0pt
\usepackage{amssymb}
\usepackage{float}
\usepackage{mathtools}
\usepackage{enumerate}
\usepackage{tikz}
\newtheorem{lemma}{Lemma}
\usepackage{cleveref}
\crefname{lemma}{Lemma}{Lemmas}
\newcommand{\append}[2]{#1^{(#2)}}
\newcommand{\fract}[1]{\left\{ #1 \right\}}
\newcommand{\ndelta}{n^{(a)}}
\renewcommand{\phi}{\varphi}
\title{State Complexity of Shifts of the Fibonacci Word}
\author{Delaram Moradi, Pierre Popoli, Jeffrey Shallit, Ingrid Vukusic}
\date{Source: arXiv:2603.18858v1 (CC BY 4.0)}
\begin{document}
\maketitle

\noindent\emph{Source and licence.} State Complexity of Shifts of the Fibonacci Word. Version arXiv:2603.18858v1 (CC BY 4.0), distributed by arXiv under the Creative Commons Attribution 4.0 International licence, http://creativecommons.org/licenses/by/4.0/, as recorded in the arXiv licence metadata for this exact version and shown on the abstract page https://arxiv.org/abs/2603.18858v1. Full licence notice: \texttt{licenses/arxiv-2603.18858-CC-BY-4.0.txt}.\par\smallskip

\noindent\textit{Editorial scope.} Counted unit: the article's lemma lem:flip\_append (source0.tex lines 1297--1308). The article's complete original proof of it is delivered with it (source0.tex lines 1311--1425, one \texttt{proof} environment of 115 lines). No mathematics is written, repaired or re-derived: the statement and the proof are the article's own lines, in the article's own notation, and no retained line is substituted or re-flowed.  No further source-local context is needed: every cross-reference of every retained line resolves inside the two delivered spans.  Nothing was omitted from any retained span, so the package carries no omission record.  No retained line contains a citation, so the wrapper carries no bibliography and no bibliography entry.  No retained line contains a figure environment, an image-inclusion command or drawing code, so no image is delivered and the package declares no asset dependency.  Editorial formatting only, in addition: an AMS \texttt{amsart} preamble that restates the article's own declarations and macros used by the retained lines, namely the environments \texttt{lemma} on separate counters, together with the standard packages those lines need.  The retained environments print in this document as Lemma 1 in order of appearance, whereas the article prints the same items with the numbering of its own full document.  The complete native source of the article as distributed in the arXiv e-print for this exact version is bundled unchanged as \texttt{sources/196730/source0.tex}.
\begin{lemma}\label{lem:flip_append}
Let $x,y$ be negative integers such that the interval $I = (x \phi,y \phi) \bmod 1$ does not contain the two points $\fract{-\phi}$ and $\fract{-2\phi}$.
Let $n$ be a nonnegative integer. If $f(n) = 1$ set $a=0$, otherwise $a \in \{0,1\}$.
Then we have
\begin{equation}\label{eq:flip_append}
    \fract{n\phi} \in I
    \implies
    \fract{\ndelta\phi} \in (\append{y}{a} \phi, \append{x}{a} \phi)\bmod 1
    =: I_a.
\end{equation}
Moreover, if $x\neq -1$, the interval $I_a$ does not contain $\fract{-\phi}, \fract{-2\phi}$.
\end{lemma}

\begin{proof}%[Proof of \cref{lem:flip_append}]
First we define the following {\tt Walnut} commands
\begin{verbatim}
reg lshift {0,1} {0,1} "([0,0]|[0,1][1,1]*[1,0])*":
# accepts (x,y) iff y is the left-shift of x.

reg rshift {0,1} {0,1} "([0,0]|[1,0][1,1]*[0,1])*(()|[1,0][1,1]*)":
# accepts (x,y) if y is the right-shift of x

def phin "?msd_fib (x=0 & y=0) | Ez $lshift(x-1,z) & y=z+1":
# (x,y) is accepted iff floor(x*phi) = y.

reg isfib msd_fib "0*10*":
# accepts x iff x is a positive Fibonacci number.

def smallest_fib_geq "$isfib(y) & y>=x & Az ($isfib(z) & z>=x) => z>=y":
# accepts (x,y) iff y is the smallest Fibonacci number >= x, for x>=1.
# 5 states

def fib_neg_repr "?msd_fib (x=0 & y=0 & z=0) | 
    (x>=1 & $smallest_fib_geq(x,y) & z+x=y)":
# accepts (x,y,z) iff -y+z is the negative Fibonacci expansion of -x,
    where y is a Fibonacci number.
# 9 states

def prime_inv_0 "?msd_fib Er,s,t,u $fib_neg_repr(x,r,s) &
    $rshift(t,r) & $rshift(u,s) & F[t]=@0 & F[u]=@0 & t=y+u":
# accepts (x,y) if -y is the integer obtained from 
    appending 0 to the representation of -x
# 7 states

def prime_inv_1 "?msd_fib (x=1 & y=1) | (x=2 & y=2) |
    Er,s,t,u $fib_neg_repr(x,r,s) &
    $rshift(t,r) & $rshift(u,s) & F[t]=@0 & F[u]=@1 & t=y+u":
# accepts (x,y) if -y is the integer obtained from 
    appending 1 to the representation of -x
# 9 states
\end{verbatim}

Note that the special cases $(x,y)=(1,1)$ and $(x,y) = (2,2)$ in {\tt prime\_inv\_1} account for the
fact that if we append $1$ to the representation of $-1, -2$, we get invalid representations. But these cases are still covered in the lemma.

Next, we define the two following functions about fractional parts. Notice that we are not using circular interpretation here. 

\begin{verbatim}
def compare_frac_part "?msd_fib 
    (x>y & Et,u,v $phin(x,t) & $phin(y,u) & $phin(x-y,v) & v+u<t) | 
    (x<y & Et,u,v $phin(x,t) & $phin(y,u) & $phin(y-x,v) & u<t+1+v)":
# accepts (x,y) iff {x*phi} < {y*phi}  for x, y >= 0 
# 22 states

def compare_frac_part2 "?msd_fib Et,u,v 
    $phin(x,t) & $phin(y,u) & $phin(x+y,v) & v>u+t":
# accepts (x,y) iff {-x*phi} < {y*phi} for x, y >= 1  
# 26 states
\end{verbatim}

\begin{verbatim}
def lies_in "?msd_fib  
 ($compare_frac_part(z,x) & $compare_frac_part2(x,y) 
             & ~$compare_frac_part2(y,z)) |
 ($compare_frac_part(x,z) & 
    ($compare_frac_part2(x,y) | ~$compare_frac_part2(z,y)))":
# checks if  {y*phi} is in the interval ({-x*phi}, {-z*phi}) 
    for the circular interpretation of intervals
# 97 states

def cw "?msd_fib 
    ($compare_frac_part(x,z) & $compare_frac_part(x,y) & 
        $compare_frac_part(y,z)) |
    ($compare_frac_part(z,x) & ($compare_frac_part(x,y) |   
        $compare_frac_part(y,z)))":
# checks if {y*phi} is in the interval ({x*phi}, {z*phi}) 
    for circular interpretation of intervals; cw = clockwise
# 44 states
\end{verbatim}

Now that we have these, we can rewrite the theorem statement in first-order logic and ask {\tt Walnut} to prove it:
\begin{verbatim}
eval check_flip_append_0 "?msd_fib An,n0,a,b,a0,b0 
    (a>=1 & b>=1 & F[n0]=@0 & 
    $rshift(n0,n) & $prime_inv_0(a,a0) & $prime_inv_0(b,b0) & 
    ($cw(a,1,b)|a=1|b=1) & ($cw(a,2,b)|a=2|b=2) 
    & $lies_in(a,n,b)) 
    => ($lies_in(b0,n0,a0))":
\end{verbatim}

\begin{verbatim}
eval check_flip_append_0_interval "?msd_fib An,n0,a,b,a0,b0 
    (a>=2 & b>=1 & F[n0]=@0 & 
    $rshift(n0,n) & $prime_inv_0(a,a0) & $prime_inv_0(b,b0) & 
    ($cw(a,1,b)|a=1|b=1) & ($cw(a,2,b)|a=2|b=2) & 
    $lies_in(a,n,b)) 
    => (($cw(b0,1,a0)|a0=1|b0=1) & ($cw(b0,2,a0)|b0=2|a0=2) )":
\end{verbatim}

\begin{verbatim}
eval check_flip_append_1 "?msd_fib An,n1,a,b,a1,b1 
    (a>=1 & b>=1 & F[n1]=@1 & 
    $rshift(n1,n) & $prime_inv_1(a,a1) & $prime_inv_1(b,b1) & 
    ($cw(a,1,b)|a=1|b=1) & ($cw(a,2,b)|a=2|b=2) & 
    $lies_in(a,n,b)) 
    => ($lies_in(b1,n1,a1))":
\end{verbatim}

\begin{verbatim}
eval check_flip_append_1_interval "?msd_fib An,n1,a,b,a1,b1 
    (a>=2 & b>=1 & F[n1]=@1 & 
    $rshift(n1,n) & $prime_inv_1(a,a1) & $prime_inv_1(b,b1) & 
    ($cw(a,1,b)|a=1|b=1) & ($cw(a,2,b)|a=2|b=2) & 
    $lies_in(a,n,b)) 
    => (($cw(b1,1,a1)|a1=1|b1=1) & ($cw(b1,2,a1)|b1=2|a1=2))":
\end{verbatim}
All of these return {\tt TRUE}.
\end{proof}

\end{document}
